% Historical detail archive; current main text controls interpretation.
% Comparator execution settings and runtime mappings have not been recovered.
\begin{figure*}[t]
\centering
\includegraphics[width=0.88\textwidth]{validation-paired-outcomes.pdf}
\caption{Paired outcomes and recorded trigger burden. (a) Full-minus-comparator
differences with 95\% participant-bootstrap intervals; sample sizes appear beside
each endpoint. (b) Participant rates over 29.164 reference-negative online hours.
Lines connect the same participant; coincident observations are retained.
Rates describe the enumerated event set.}
\label{fig:participant_outcomes}
\end{figure*}

\subsection{Scripted Targets and Natural-Use Triggering}
\label{sec:participant_validation}
The study recorded 42 sessions from 30 participants between 17 August and
5 September 2026: 12 natural-use, six scripted-action, seven simulated-positive,
five challenge, and 12 repeat sessions. The recorded release,
v1.3.0+frozen, used a personal Balance baseline and quality gates (Full).
Two historical comparison labels shared all 201 event IDs: a population Balance reference
replacing the personal baseline (Population baseline), and disabled quality
gates (No gates). Reference annotations were recorded as blinded to system
outputs. The 603 scheme-specific records were analysed by event, with
participant-cluster percentile bootstrap intervals from 10,000 paired resamples
(seed 20260910); rates are event-weighted and intervals are pointwise.
Trigger detection denotes an event trigger. Its timestamp is the first trigger
for detected events and the post-window negative decision for untriggered events.
Completeness denotes process completion linked to the final report, independently
of trigger status. These timestamps are not pooled as completed-check latency.



Full recorded triggers for 35/38 scripted targets from 12 participants (92.1\%), compared
with 29/38 for Population baseline and 33/38 for No gates.
Its component counts were 9/10 Face, 7/8 Arms, 10/11 Balance, and 9/9 Speech.
The recorded 15.8-point trigger-detection difference from Population baseline (95\% interval
7.9--23.7) came entirely from six additional Balance detections.
The difference from No gates was 5.3 points (0.0--14.7), retaining uncertainty
about detection improvement (Fig.~\ref{fig:participant_outcomes}a).
Among 18 reference-incomplete events from 13 participants, including two
aborted attempts, incorrect completions were 1/18, 1/18, and 15/18, respectively.
The Full-minus-No gates difference was $-77.8$ points ($-95.0$ to $-55.6$);
the remaining Full error involved non-target speech.
All three configurations marked all 159 reference-complete inputs as complete.
Thus Full retained these inputs while recording fewer incorrect completions
than No gates under the engineering reference protocol. The 177 started
measurements involved 25 participants; 159 were reference-complete (89.8\%).



Across 12 natural-use participants, the enumerated events contained 14/20/19
false triggers for Full/Population baseline/No gates over 29.164 reference-negative
online hours (0.480/0.686/0.651 per hour; Fig.~\ref{fig:participant_outcomes}b).
These are recorded triggers, not a measured rate of final stroke alerts.
Continuous-input replay and linked guided-check outcomes are needed to establish
complete trigger coverage and subsequent alert burden. Online availability was 29.522/31.810 scheduled hours
(92.8\%). Outputs below the event-specific reference urgency occurred in
3/38, 9/38, and 5/38 targets, respectively, all among missed targets.
Four targets had non-new-onset labels requiring \textit{warning}; the other
34 required \textit{urgent}. Among detected targets, urgency-label mismatches
were 0/35, 0/29, and 0/33. This conditional result excludes detection failures.
Symptom counts were below reference in 2/35, 3/29, and 6/33 targets with known
counts; the remaining 3, 9, and 5 counts were unknown. Correction histories
were intact in 8/9, 8/9, and 6/9 cases. These record-level failures do not
identify whether submission, state handling, export, or persistence failed.

Figure~\ref{fig:participant_reference} shows exported-value differences and
within-participant repeat measurements. Face mouth-corner MAE was 0.00452
interocular-width units (33 measurements, 22 participants).
Eleven complete Face repeat pairs yielded ICC(2,1)=0.935 (95\% interval
0.801--0.975), with measurements 11.2--15.6 days apart.
One of 12 repeat participants lacked the first corresponding Face measurement.
For the exported Eyes gaze-offset field, MAE was 0.280$^\circ$ against desktop eye-tracker values
(25 measurements, 19 participants), over system values of 0.094--1.641$^\circ$.
The Eyes field is not mapped to a unique runtime monocular or binocular output;
this numerical comparison does not establish runtime angular accuracy.
Automatic Eyes abnormality classification remained disabled.

\begin{figure*}[t]
\centering
\includegraphics[width=0.88\textwidth]{validation-reference-repeatability.pdf}
\caption{Exported-value comparisons and Face repeatability.
(a,b) System-minus-reference errors for Face mouth-corner change and Eyes gaze
offset; dashed lines indicate zero difference. Eyes compares exported fields with
unresolved runtime mapping. (c) Two system measurements for each
of 11 participants with complete Face repeat pairs. Lines connect participants
across visits.}
\label{fig:participant_reference}
\end{figure*}

The recorded protocol specifies a 5-s Arms hold and a 2-s neutral plus
3-s smile Face sequence. Eyes totals include settling-period attempts,
which runtime trial counters exclude. The population-reference parameters,
its enrollment provenance, and the exact No gates rule changes are not specified
in the available configuration record. Comparisons therefore describe these
recorded configurations without isolating the effects of personalization or
individual gates. Scripted targets support engineering evaluation, not
clinical sensitivity.
